% ==============================================================================
% CHƯƠNG 8: KẾT LUẬN
% ==============================================================================
\chapter{Kết luận}
\label{chap:ket-luan}

\section{Kết quả đạt được}
\label{sec:kl-ket-qua}

\begin{dsdung}
  \item Cài đặt đủ năm hàm \ma{findMeshPath}, \ma{findAUVPath}, \ma{findProbabilisticPath},
        \ma{findGamePath}, \ma{findPCBPath} đúng chữ ký của đề, chỉ dùng sáu thư viện được phép,
        biên dịch sạch ở C++11 đến C++20.
  \item Quy cả năm bài toán về một đồ thị có hướng có trọng số không âm và một lõi A* dùng chung;
        phí bẻ góc của Task~5 được xử lý bằng đồ thị trạng thái $(\text{ô}, \text{hướng})$.
  \item Chứng minh tính nhất quán của heuristic số bước $\times c_{\min}$ (Task~1), Chebyshev
        (Task~4), Manhattan (Task~5), và điều kiện $\alpha \ge 1$ cho Euclid/Chebyshev (Task~2).
  \item Trả lời các yêu cầu báo cáo bằng vết thực thi và số liệu sinh tự động, có thể tái lập.
\end{dsdung}

\begin{bangbieu}{Tổng hợp heuristic theo từng task}{tab:kl-tong-hop}{C{1cm} L{3.4cm} Y}
  \hangtieude \thead{Task} & \thead{Heuristic} & \thead{Tính chất} \\
  \midrule
  1 & $\hops(v) \cdot c_{\min}$ & \otot Nhất quán, luôn tối ưu \\
  2 & Manhattan / Euclid / Chebyshev 3D & Euclid, Chebyshev nhất quán khi $\alpha \ge 1$; Manhattan có thể lệch \\
  3 & Manhattan / Euclid 2D & Euclid nhất quán khi $d(u, v) \ge$ khoảng cách thẳng, với mọi $\lambda$ \\
  4 & Manhattan / Chebyshev & \otot Chebyshev nhất quán; \oxau Manhattan đánh giá quá do bước chéo \\
  5 & Manhattan / Chebyshev & \otot Cả hai nhất quán; Manhattan mở rộng ít hơn \\
\end{bangbieu}

\section{Hạn chế và hướng phát triển}
\label{sec:kl-han-che}

\begin{dssai}
  \item Quy tắc phá hòa ($f$, rồi $h$, rồi thứ tự vào \OPEN) và cách đặt tên nút (chỉ số, hướng,
        tọa độ) là lựa chọn hợp lý khi đề chưa quy định; cần đối chiếu lại với testcase LMS khi
        cổng nộp bài mở.
  \item Ở Task~2 và Task~3, heuristic theo đúng công thức của đề có thể đánh giá quá khi đơn vị chi
        phí không gắn với tọa độ, nên không bảo đảm tối ưu.
\end{dssai}

Một số cải tiến có thể làm thêm: nhân heuristic với hệ số $\alpha$ của \eqref{eq:t2-alpha} để
khôi phục tính chấp nhận được; dùng chi phí $-\ln(1 - P_{\text{risk}})$ nếu muốn tối đa hóa chính
xác $P_{\text{safe}}$ ở Task~3; cộng thêm $B$ vào heuristic của Task~5 khi ô ra không nằm trên tia thẳng
theo hướng hiện tại (ở điểm xuất phát: không cùng hàng, cùng cột), vì khi đó chắc chắn còn ít nhất
một lần bẻ góc.
